\section{Limitations}
\label{sec:limitations}

\paragraph{A positional shortcut in the task.} Zoology's MQAR places every key in the first 256 tokens, so the
coarse level can be solved by preferring early super-blocks (\S\ref{sec:selection}). Conclusions about $S_1$ may not
transfer to tasks whose keys are spread across the context. We built a variant with keys placed anywhere, but dense
attention did not learn it under this recipe (recall 0.0002--0.0044, near chance, across four variants), so it could not serve
as a test. A learnable shortcut-free task is the most important next step.

\paragraph{One valid dense comparison.} Dense attention transitioned on only one of the three reserved seeds, so the
gaps of 0.291 and 0.055 rest on a single pair. They are observations, not estimates, and we make no claim of
statistical significance. The dense column also mixes numerics: the valid pair is TF32, while the float32 reruns are
diagnostics.

\paragraph{A selected operating point.} $S_1{=}6$ was chosen after the pilot because it scored highest there. The
reserved seeds confirm that point only; $S_1{=}4$ and $S_1{=}12$ were not run on them.

\paragraph{Training ceiling and a stopped run.} A non-transition means only that validation recall did not exceed
0.30 within 60 epochs of a cosine schedule that ends near zero learning rate; a longer or different schedule could
change any such run. The recipe was not tuned for dense attention under the M5 profile. One float32 dense
run (seed 810001) was stopped by us after 55 of 60 epochs with validation recall at 0.0068, and has no held-out score.

\paragraph{Scale and scope.} Models have 2 layers and width 128; there is one task and one length (2048). We make no
claim about real text, longer contexts or larger models. At length 4096 even dense attention failed under this
recipe, and longer lengths fall outside the task's vocabulary constraint.

\paragraph{Mixers are matched in width, not in everything.} Dense attention uses dropout 0.1 and the selector
none; the selector has extra gate parameters. Either could contribute to the difference in training reliability.

\paragraph{The budget knob is not a pure intervention.} Raising $S_1$ also enlarges the set of blocks the fine level
chooses from, and on this task the parent level can be solved by position, so the budget curve mainly reflects how
reliably the fine choice is learned (\S\ref{sec:budget}). Flat block selection was run only as $S_1{=}12$, on one seed
per profile, so whether the hierarchy helps or costs recall is not settled. The fine budget $K$ was fixed at 4; a
post-hoc probe suggests it matters more than $S_1$ (\S\ref{sec:selection}).

\paragraph{No speed claim.} Selector scores per query is a count, not runtime. Without a fused kernel our selector
was slower than dense attention at every budget (Table~\ref{tab:cost}).

\paragraph{Mechanism evidence is post hoc.} The selection diagnostic (\S\ref{sec:selection}) replays trained models
and describes their choices; it does not intervene on them. The repair screen (\S\ref{sec:fixes}), the distance
analysis and the budget probe use one M4-profile checkpoint per arm, and their explanations account for the numbers
arithmetically rather than causally.
